\section{Removing the small-order rational characters}
\label{sec:kt-small-character-order}

The low-degree real class-number theorems give a direct restriction
on the rational components still requiring the unit-trace method.
The class-number theorems in this section are known inputs, not new
results. The short argument records their precise implication for
the character and integral-cohomology conventions of this report.

\begin{proposition}[KT-SMALL-ORDER-COMPONENT-REDUCTION]
Let $p>3$ be prime and let $e\ge2$ be even. If the order of the
Teichm\"uller character $\omega^e$ is in $\{1,2,3,4\}$, then
\[
                 A_p^{[-e]}=0,
                 \qquad K_{2e}(\mathbf Z)\{p\}=0.       \tag{LO.1}
\]
Consequently a remaining rational component in the all-weight
problem has $\operatorname{ord}(\omega^e)\ge5$.
No bound on $e$ relative to $p$ is needed.
\end{proposition}

\begin{proof}
Put $L=\mathbf Q(\mu_p)$ and let $d$ be the order of
$\omega^{-e}$. Let $F=L^{\ker(\omega^{-e})}$, so
$[F:\mathbf Q]=d$. Since $e$ is even, $F$ is totally real.
We first recall the exact class-number inputs for these four
possible degrees.

For $d=1$, $F=\mathbf Q$ has class number one. For $d=2$,
$p\equiv1\pmod4$ and $F=\mathbf Q(\sqrt p)$; the real
quadratic bound in Le, Theorem~(a), printed p.~141, already
checked in KT-ODD-DISCRIMINANT-TWO-UNIT-FAMILY, gives
$h(F)\le\lfloor\sqrt p/2\rfloor<p$.

For $d=3$, $F$ is the unique cyclic cubic field of conductor
$p$. Moser--Payan, \emph{J. Math. Soc. Japan}~33 (1981),
701--706, \S4, unnumbered Theorem and Corollary on p.~705,
gives $h(F)<2p/3$, hence $p\nmid h(F)$.

For $d=4$, evenness of the character implies
$p\equiv1\pmod8$. Put $k=\mathbf Q(\sqrt p)\subset F$.
Moser, Grenoble seminar 1979--1980, expos\'e~3,
Theorem~B on printed p.~III.2, gives
\[
            h(k)\le\sqrt p/2,\qquad h(F)/h(k)\le p/4.  \tag{LO.2}
\]
The same expos\'e, Theorem~1(2) on p.~III.1, gives
$h(k)\mid h(F)$: the cyclic quadratic extension $F/k$
has exactly one ramified prime. Thus both $h(k)$ and
$h(F)/h(k)$ are positive integers less than $p$, and
$p\nmid h(F)$. These are applications of the cited known
theorems; their analytic proofs are not being reproved here.

In every case $\operatorname{Cl}(F)\{p\}=0$.
Write $H=\operatorname{Gal}(L/F)$, whose order is prime to
$p$. Let $j$ and $N$ be extension and norm on ideal class
groups. They satisfy
\[
                      jN(x)=\sum_{\sigma\in H}\sigma x.
                                                               \tag{LO.3}
\]
If $x\in A_p^H$, then $N(x)$ is a $p$-primary class of $F$,
so it is zero. Formula (LO.3) gives $|H|x=0$.
Multiplication by $|H|$ is invertible on the finite $p$-group
$A_p$, hence $A_p^H=0$. The character $\omega^{-e}$ is
trivial on $H$, so $A_p^{[-e]}\subseteq A_p^H=0$.
The carried comparison K17-(K2), together with its periodic
character convention, gives the second assertion of (LO.1).
\end{proof}

\paragraph{Scope and source boundary.}
The two primary sources are
\url{https://doi.org/10.2969/jmsj/03340701} and
\url{https://www.numdam.org/item/STNG_1979-1980__8__A3_0/}.
The cubic numerical bound used is the published $2p/3$ bound;
no stronger bound is needed. The exact quartic statements and
their hypotheses were checked against the scanned theorem page.
The supplied K17 comparison is cited as settled input.
This restriction removes the specified character orders from
the rational $H^2_{\acute et}(\operatorname{Spec}\mathbf Z[1/p],
\mathbf Z_p(e+1))$ problem in (FW.1). It neither computes every
class component of higher order nor proves the all-weight
vanishing theorem. In particular no degree-six assertion is
imported from a title or a secondary bibliographic description.
